\chapter{端到端建模流程与技术原理}

本章按程序真正的数据依赖顺序说明算法。为便于交接，每一节均回答“目的、核心方法、设计
理由、输出和主要失败条件”，公式只保留理解实现所必需的部分。

\section{步骤一：病例加载、表面采样与导管识别}

\subsection{病例加载与公共几何}

\texttt{analyze\_case()} 清空 Blender 场景后导入导板、牙列和一个或多个导管装配体。
导板和牙列表面被采样为带法向和面索引的 \texttt{SurfaceSample}。导板中心优先使用封闭
网格的体积重心；如果计算出的重心落在包围盒外，则退化为表面样本均值。

导板局部法向由表面点云主平面确定。若点云中心为 $\bar{\mathbf p}$，协方差矩阵为
\[
  \mathbf C=\frac{1}{N}\sum_i
  (\mathbf p_i-\bar{\mathbf p})(\mathbf p_i-\bar{\mathbf p})^\mathsf T,
\]
最小特征值对应的特征向量近似为导板法向。程序再用上下颌或 YAML 中确认的牙合轴统一
法向符号。

\subsection{为什么先检查真实轴孔}

导管装配体可能包含导管、连接附件和实心 cutter。如果只按“长度和直径最像导管”打分，
实心圆柱也可能获胜。当前实现先对每个连通分量拟合轴线，并在轴向 15\%、25\%、35\%、
50\%、65\%、75\%、85\% 七处取探针点。点位于实体外部表示该处处在孔腔中，七点的通过
比例必须至少为 0.60。

只有通过孔道门槛的候选才会两两组合。候选对评分主要考虑：
\[
  E_{\mathrm{size}}=
  \sum_{j=1}^{2}\left(
  \frac{|H_j-H_0|}{H_0}+\frac{|R_j-R_0|}{R_0}\right),
  \qquad
  E_{\mathrm{axis}}=1-|\mathbf a_1\cdot\mathbf a_2|.
\]
代码按尺寸误差、轴线不平行误差、导管间距和组件索引依次排序，选择最优对。尺寸只在
已经确认有孔的候选中排序，不能让实心 cutter 反超。

\subsection{轴线与 C 口方向}

导管轴线不是只对全部顶点做一次 PCA。\projectpath{sleeve_estimation} 会对网格切片，
从内孔截面拟合圆和直线，以孔道证据修正主轴。轴向上下端再由顶点投影范围和端部拓扑
确定。两根导管中心连线投影到各自轴线法平面：
\[
  \mathbf c_i=
  \operatorname{normalize}\!\left[
  (\mathbf o_j-\mathbf o_i)-
  \mathbf a_i\big((\mathbf o_j-\mathbf o_i)\cdot\mathbf a_i\big)
  \right],
\]
所得 $\mathbf c_i$ 指向另一根导管，作为 C 口方向；其反方向用于放置主梁接触点。

\subsection{input 与 generated 两种导管模式}

\begin{table}[H]
\centering
\small
\begin{tabularx}{\textwidth}{p{0.18\textwidth}XX}
  \toprule
  对比项 & generated & input\\
  \midrule
  最终导管来源 & 按识别位姿和八个解析尺寸重建 & 保留识别出的输入导管分量\\
  锚点表面 & 参数化外壁 & 真实输入外壁射线命中\\
  导孔处理 & 全体融合后权威复切 & 导管加入前只复切基础体\\
  低梁嵌入 & 可全直径预埋，依赖后续复切 & 最多约 1.45 mm，并受真实壁厚限制\\
  连通体清理 & 复切后保留最大连通体 & 只删除小于 0.25 个体素的数值碎片\\
  主要优势 & 尺寸一致、对扫描噪声不敏感 & 保留真实导管形态和既有孔腔\\
  主要风险 & 重建参数与实际零件可能不一致 & 输入缺陷被保留，体素融合仍会离散化表面\\
  \bottomrule
\end{tabularx}
\caption{两种导管实体模式的布尔顺序差异。}
\end{table}

两种模式共享同一套识别位姿、C 口语义和下游规划。这个设计避免为不同导管来源复制整套
流程，但也要求每个涉及孔道的步骤明确知道当前模式。

\section{步骤二：牙位识别、FDI 编号与导板映射}

\subsection{从三维牙列到二维牙弓证据}

TwinGuide 通过工作区中的 \projectpath{tooth_guide_mapping} 模块执行牙位工作流。流程先把
牙冠表面投影到患者左右—前后平面，组合四类证据：连续三角形轮廓、最高表面高度、插值
法向和融合解剖边缘。图 \ref{fig:projection} 展示了 tooth-47 的实际输出。

\begin{figure}[H]
  \centering
  \includegraphics[width=0.96\textwidth]{04_multichannel_projection.png}
  \caption{tooth-47 多通道牙冠投影。高度、法向和边缘证据共同用于识别相邻牙冠边界。}
  \label{fig:projection}
\end{figure}

当前 tooth-47 识别配置采用 \texttt{arch\_progress} 分组策略。接触弦轮廓用于把相邻牙冠
分开，批准后的闭合轮廓面积质心成为最终牙中心，而不是直接沿用较早的基础映射中心。

\begin{figure}[H]
  \centering
  \includegraphics[width=0.82\textwidth]{05_contact_chord_contours.png}
  \caption{批准后的 present-only 接触弦轮廓。圆点是面积质心，叉号是内部参考中心。}
  \label{fig:contact-contours}
\end{figure}

\subsection{沿牙弓排序和 FDI 语义}

各牙中心沿一条有方向的牙弓曲线投影为弧长坐标 $s$。映射报告保存每颗牙的 FDI、
\texttt{arch\_s\_mm}、局部切向、局部外向、牙冠点和导板顶部点。缺失牙和排除牙保留在
FDI 顺序中用于语义推断，但不会创建虚假的“现存牙表面中心”。

\begin{figure}[H]
  \centering
  \includegraphics[width=0.98\textwidth]{06_tooth_guide_mapping.png}
  \caption{牙中心、定向牙弓和导板物理覆盖映射。右图把每颗牙的轮廓区间投影到弧长 $s$。}
  \label{fig:tooth-mapping}
\end{figure}

\subsection{左右方向安全门}

自动坐标可能存在左右翻转。若 YAML 没有直接给出确认轴，程序会比较两个候选方向，并在
存在缺失种植牙位时，比较预测缺牙位置到导管手术位的牙弓平面距离。最佳方向通常必须不
超过 15 mm，且比另一方向至少好 3 mm；证据不足即停止。tooth-47 的 YAML 已明确给出
$+X$ 为患者右到左、$+Y$ 为前到后、$+Z$ 为牙合方向，因此报告采用确认轴而非自动二选一。

\subsection{缓存与防陈旧结果}

映射清单记录病例 YAML、牙列、导板和导管参考 STL 的 SHA-256。只有算法版本、输入路径和
摘要全部一致时才允许复用，否则重新执行。生产流程还要求映射状态完整、全部 QA 为真、
上下颌一致、报告中的输入路径与当前配置一致。

\section{步骤三：导孔、操作窗与观察窗}

\subsection{导孔}

每根导管生成一个有限圆柱 cutter。若导管中心为 $\mathbf o$、轴为 $\mathbf a$、轴向范围
为 $[t_{\min},t_{\max}]$、余量为 $m$，则 cutter 两端为
\[
  \mathbf p_0=\mathbf o+(t_{\min}-m)\mathbf a,
  \qquad
  \mathbf p_1=\mathbf o+(t_{\max}+m)\mathbf a.
\]
tooth-47 的轴向余量为 5 mm，规划孔径为 2.10 mm。有限而非无限圆柱能控制差集范围，
又通过上下余量确保穿透融合后的实体。

\subsection{双导管操作窗}

程序在非导管分量中寻找靠近双导管中点、主平面长短比不大于约 1.3 的紧凑圆形操作结构。
操作窗局部标架由两导管平均轴、中心连线投影和二者叉积构成。长边覆盖导管中心距、两侧
导管半径及切向余量；短边覆盖操作特征直径和双侧 bitangent 余量；深度取局部导板厚度与
导管长度的较大者，再加轴向余量。最终 cutter 是带圆角的定向长方体。

\subsection{按 FDI 范围构造轴扫掠观察窗}

观察窗不再使用固定世界坐标下的粗矩形。YAML 指定起止 FDI，牙位映射给出牙弓语义轴和
朝唇颊侧外部的方向。程序沿语义轴布置截面，在每个截面内从牙合方向旋转到外侧方向，形成
约 90° 的扇形扫掠体；自适应射线测量真实导板壁厚，并附加配置的 overcut 以保证穿透。

\begin{figure}[H]
  \centering
  \includegraphics[width=0.96\textwidth]{07_observation_windows.png}
  \caption{tooth-47 两个观察窗的三向投影。蓝线对应 42--32，绿线对应 34--35。}
  \label{fig:observation}
\end{figure}

观察窗生成器检查：组合 cutter 是否为闭合体、是否确实移除材料、语义轴是否连续开放、
是否能看到牙面、是否残留 cutter 重叠、差集体积恒等式是否成立。某些轴行失败时可以按配置
逐级增加局部 drop；修改只写入派生映射，不回写原牙位报告。

\begin{figure}[H]
  \centering
  \includegraphics[width=0.78\textwidth]{08_cutouts.png}
  \caption{导孔、操作窗和观察窗 cutter 与导板的空间关系。透明棕色体为操作区，青色体为 FDI 观察窗。}
  \label{fig:cutouts}
\end{figure}

\section{步骤四：导管 Q/P 与导板 A 锚点}

\subsection{导管表面接触点 Q 和中心线点 P}

每根导管在 C 口反方向的稳定外壁母线上选择上下两个截面。高端位置和低端位置分别近似为
\[
  t_{\mathrm{upper}}=t_{\max}-2.0-r,
  \qquad
  t_{\mathrm{lower}}=t_{\min}+1.0+r,
\]
其中 $r$ 是梁半径。由截面中心向外射线命中导管真实外壁得到 Q。若期望梁与导管重叠深度
为 $d$，外壁法向为 $\mathbf n$，中心线点为
\[
  \boxed{\mathbf P=\mathbf Q+(r-d)\mathbf n}.
\]
这一区分很重要：Q 描述实体表面接触，P 描述圆管中心线。直接把中心线放在 Q 上只能得到
半径级接触，不能表达受控预埋。

\subsection{牙位参考点 T 与旋转射线}

导板端不再按整条切面轨迹的 20\%/80\% 弧长取点。对单牙中心或双牙中点站位，先取附近牙冠
最高点 $\mathbf H$，再沿正牙合方向向内下移 2 mm：
\[
  \mathbf T=\mathbf H-2\,\mathbf e_{\mathrm{occ}}.
\]
单牙站位使用最近且被导板覆盖的邻牙建立旋转面；双牙站位使用两牙中心连线。旋转射线为
\[
  \mathbf d(\theta)=
  \cos\theta\,\mathbf e_{\mathrm{occ}}+
  \sin\theta\,\mathbf e_{\mathrm{side}},
\]
其中 $\mathbf e_{\mathrm{side}}$ 的符号由该牙位的 \texttt{local\_outward} 决定，因此 U 侧和
背 U 侧不是屏幕左/右。主梁常用 U 侧 70°、背 U 侧 90°，但每个锚点都可独立配置。

\subsection{为什么选择“外壁出口”}

射线与薄壳导板可能产生入口、内壁、外壁和重复三角交点。程序把射线参数距离相差小于
0.02 mm 的交点聚类，再检查命中面法向与射线方向的点积；首个法向朝射线方向且点积不小于
0.05 的交点才被视为外壁出口 A。这样可以避免把内壁入口误当作梁端。

\subsection{独立端点与导管配对}

每个端部配置一个 U 侧和一个背 U 侧锚点。程序比较直接配对与反向配对的总距离，使每根
导管从两个端部各获得一个锚点。这保证每根梁都跨越两个导板区域，而不是让一根导管只连接
到同一局部区域。

\begin{figure}[H]
  \centering
  \includegraphics[width=0.82\textwidth]{09_link_points.png}
  \caption{tooth-47 锚点可视化。黄色标记为导管侧 Q/P 相关点，棕色标记为导板端锚点。}
  \label{fig:anchors}
\end{figure}

\section{步骤五与六：四根主梁、Y 型按压梁和实体扫掠}

\subsection{主梁拓扑}

每根导管生成高、低两根连续梁，均从一个导板端出发，经导管 P 点到另一个导板端。因此一个
种植位的两根导管共四根主梁。高梁直接经过高端 P；低梁先构造外层代理曲线，再只在 P 附近
局部下潜，降低大范围弯折和孔道侵入风险。

\subsection{五次 Hermite 曲线}

每段曲线采用端点二阶导数为零的五次 Hermite 基函数：
\begin{align*}
  \mathbf C(u)=&\;h_0(u)\mathbf p_0+h_1(u)\mathbf m_0
  +h_2(u)\mathbf p_1+h_3(u)\mathbf m_1,\quad u\in[0,1],\\
  h_0=&1-10u^3+15u^4-6u^5,\\
  h_1=&u-6u^3+8u^4-3u^5,\\
  h_2=&10u^3-15u^4+6u^5,\\
  h_3=&-4u^3+7u^4-3u^5.
\end{align*}
两段在导管 P 共享一阶导数，避免折线式尖角。端点切向先投影到相应表面切平面，再按端点
距离和张力系数缩放。中心线按约 0.30 mm 间距离散，且每段至少 17 个奇数采样点。

\subsection{低梁局部下潜}

低梁的外层代理接触点位于外壁 Q 附近。程序在代理曲线两侧约 5 mm 弧长处选择拼接点，
左拼接继承左侧自然切向，右拼接继承右侧自然切向，只有深埋 P 使用左右拼接点连线方向。
这一设计修复了“整段共用一个方向”导致的两侧切向不连续。

\subsection{Y 型按压梁}

当前支持三种模式：

\begin{itemize}
  \item \texttt{inner\_sleeve\_upper\_y}：一个锚点在内侧导管高端，两个在导板；
  \item \texttt{three\_tooth\_anchors\_y}：三个锚点全部由牙位射线落在导板；
  \item \texttt{terminal\_u\_extension\_anchor\_y}：一个锚点落在末端 U 梁，两个在导板。
\end{itemize}

三牙位模式取三个中心线锚点的算术中心，并沿病例确认的牙合轴抬高默认 2 mm。程序检查
三锚点展开面积、最短边和汇合点处最小臂角（默认至少 25°），避免近乎共线的“假 Y 型”。
末端 U 梁模式则在允许中心线上最大化到两个导板锚点的较小距离，即 maximin 选点。

\subsection{平行输运圆管扫掠}

直接使用 Frenet 标架会在曲率接近零或反曲点处突然翻转。项目先估计离散切向，再用相邻
切向的叉积确定最小旋转轴，通过 Rodrigues 旋转把前一截面法向输运到下一截面。随后每个
采样点构造 64 边圆环并连接相邻圆环，封闭首尾端盖。

\subsection{根部渐粗、球头与贴合脚}

梁在导板端最后约 3 mm 通过 smoothstep 从标准半径渐粗到约 1.08 倍；同一唯一导板端再增加
根部球和曲面贴合椭圆脚。贴合脚向导板内预埋 0.25 mm，并按局部表面投影。若完整足印投影
距离超过 1.56 mm，程序逐步缩小足印，最低到 77\%；仍超限时只发出警告并继续生成。

\begin{figure}[H]
  \centering
  \includegraphics[width=0.84\textwidth]{10_connectors.png}
  \caption{tooth-47 主梁、Y 梁、根部强化与导管的装配预览。金色为新增梁架。}
  \label{fig:connectors}
\end{figure}

\section{特殊病例拓扑}

这些结构不是 tooth-47 的主线，但属于“整个项目”必须交接的能力。

\begin{description}
  \item[两分量断裂导板跨接] 要求导板恰好有两个主要分量，在两个牙位站位分别求 U/背 U
  锚点，使同侧锚点跨越不同分量，生成两根直线跨接梁。
  \item[末端 U 型延伸] 用于末端牙存在但导板覆盖过短。程序从参考邻牙指向末端牙确定远中
  方向，优先依据末端牙闭合轮廓外扩“梁半径 + 牙体净距 + 安全余量”，构造切向连续 U 梁。
  \item[远中公共节点] 用于末端牙缺失且远中无导板。以两根末端导管低端 P 中点 $\mathbf B$
  为基准，沿同时垂直于导管连线和公共轴线的远中方向移动两个平均外径：
  $\mathbf G=\mathbf B+2D\mathbf e_{\mathrm{distal}}$，四根主梁共享 G。
  \item[相邻双种植位连续路径] U 侧和背 U 侧各生成一条跨两个种植位的连续路径，上下两层
  各执行一次，仍得到四根主梁。
\end{description}

末端 U 型延伸与远中公共节点在当前配置中互斥，因为二者代表“有末端牙可绕行”和“末端缺牙
自由闭合”两种不同策略。

\section{牙列净距、体素融合与精确布尔}

\subsection{牙列保护裁剪}

主梁、Y 梁和根部强化先合并为梁架，再用牙列网格作为 cutter，并附加配置的
\texttt{connector\_dental\_clearance\_mm}。这样导板端可以直接以表面 A 点作为中心线端点，
让梁充分进入导板；侵入牙列的部分由差集统一清除。

\subsection{为什么使用体素融合}

多个圆管、球、贴合脚、导管和扫描导板往往只有局部重叠，直接精确并集容易因共面、薄片和
非流形输入失败。\texttt{voxel\_union()} 复制对象、拼接后按 0.2 mm 体素重网格，得到一致
封闭壳体。其代价是细节误差约为体素量级，薄壁、窄孔和尖锐边会被平滑或回缩。

\subsection{差集后端与修复}

初始导板预览使用 Blender 布尔求解器按 EXACT、MANIFOLD、FLOAT 回退；关键的复切和手机
避让使用 manifold3d 连续差集。cutter 默认做 0.02 mm Minkowski 外扩并以 0.005 mm 简化，
减少共面残留。多个 cutter 在同一 manifold 会话中连续扣除，避免中间结果反复序列化为
浮点网格。

如果 STL 序列化产生亚微米重合点或不超过 0.02 mm 的短坏边，代码允许局部焊接或折叠；
修复后必须重新加载并确认为单一封闭体。超过阈值则失败，避免以大尺度自动修复掩盖真实
几何错误。

\section{手机运动包络避让}

程序读取手机 STL 的最大面积连通分量，从止挡报告取得配对止挡面中心和公共轴。两中心中点
作为枢轴 $\mathbf c$，在 $[-5^\circ,5^\circ]$ 内采样 41 个姿态。对任意手机点
$\mathbf p$，第 $k$ 个姿态为
\[
  \mathbf p_k=\mathbf c+\mathbf R(\mathbf a,\theta_k)(\mathbf p-\mathbf c).
\]
所有姿态精确并集成封闭包络，再从最终结构做差集。当前不模拟轴向下压；也不保护进入包络
的梁，因此避让后梁可能被切断，必须重新检查连续性。

\section{两种最终布尔顺序}

\begin{figure}[H]
\centering
\begin{tikzpicture}[
  node distance=5mm,
  box/.style={draw=TGBlue,rounded corners,fill=TGLightBlue,text width=6.2cm,
    minimum height=8mm,align=center,font=\small},
  arrow/.style={-{Latex[length=2mm]},thick,color=TGBlue}
]
  \node[box] (g1) {generated：先切原导板};
  \node[box,below=of g1] (g2) {导板 + 重建导管 + 梁架体素融合};
  \node[box,below=of g2] (g3) {复切导孔与观察窗};
  \node[box,below=of g3] (g4) {手机包络差集、最大连通体、修复导出};
  \draw[arrow] (g1)--(g2);\draw[arrow] (g2)--(g3);\draw[arrow] (g3)--(g4);

  \node[box,right=15mm of g1] (i1) {input：先切原导板};
  \node[box,below=of i1] (i2) {导板 + 梁架体素融合（暂不含导管）};
  \node[box,below=of i2] (i3) {基础体复切孔/窗并做手机差集};
  \node[box,below=of i3] (i4) {最后融合受保护输入导管，仅清亚体素碎片};
  \draw[arrow] (i1)--(i2);\draw[arrow] (i2)--(i3);\draw[arrow] (i3)--(i4);
\end{tikzpicture}
\caption{generated 与 input 模式必须采用不同的复切顺序。}
\label{fig:boolean-order}
\end{figure}

\section{导出、过程图与独立验证}

生成器输出最终 STL、四个标准视图以及输入、导管、锚点、梁架和切口过程图。验证器随后
检查：边界边、非流形边、重复三角形、连通体；导管表面保留率；主梁中心线包覆率；根部球
和贴合脚；Y 梁三臂；导孔多截面多半径探针；观察窗脊线净空。关键默认阈值包括导管保留率
0.90、梁中心线保留率 0.95、观察窗脊线净距 0.20 mm。

验证器的意义是检测“规划看起来正确、但最终布尔破坏了结构”的情况；它不是临床适配验证，
也没有替代实体试戴、打印工艺补偿和医疗审核。

